% ====================================================================
% ch2_preamble.tex — Chapter 2 (v1.0.21) 공통 preamble·매크로
%   graphite_ica_ch2_v1.0.21.tex 가 \input 한다. 절별 파일은 _sections/ 에.
%   빌드: xelatex -interaction=nonstopmode (kotex — 한글 전제), 3-pass.
%   렌더 목표 = v1.0.18.2 Ch2 동등(문서클래스·패키지·박스 체계·매크로 계승),
%   조립만 절별 \input 으로 전환. 재사용 불가로 Ch1 preamble 과 분리:
%     번호 prefix 2.N·박스 세트(srcbox/procedurebox)·tikz decorations·전용 매크로.
% ====================================================================
\documentclass[11pt,a4paper]{article}
\usepackage[margin=22mm]{geometry}
\usepackage{kotex}
\setmonofont{D2Coding}[Scale=MatchLowercase]
\usepackage{amsmath,amssymb,mathtools,bm}
\usepackage{booktabs,longtable,array}
\usepackage{enumitem}
\usepackage{amsthm}
\usepackage{tikz}
\usetikzlibrary{calc,arrows.meta,decorations.pathreplacing,positioning}
\usepackage{hyperref}
\usepackage{fancyhdr}
\usepackage{lastpage}
\usepackage{setspace}
\setstretch{1.12}
\setlength{\parskip}{0.45em}\setlength{\parindent}{0pt}\setlength{\headheight}{14pt}
\setlength{\emergencystretch}{2em}
\hypersetup{colorlinks=true, linkcolor=blue, urlcolor=blue, citecolor=blue,
  pdftitle={흑연 음극 dQ/dV — Chapter 2 (v1.0.21): 코인 하프셀 엔트로피·가역 발열의 통계열역학}}
\pagestyle{fancy}\fancyhf{}
\lhead{흑연 음극 $dQ/dV$ — Ch.2 (v1.0.21, 엔트로피 통계열역학)}\rhead{\thepage/\pageref{LastPage}}
\renewcommand{\headrulewidth}{0.4pt}
\newtheorem*{keybox}{핵심}
\newtheorem*{srcbox}{문헌 근거}
\newtheorem*{warnbox}{주의}
\newtheorem*{procedurebox}{계산 절차}
\newtheorem*{bgbox}{배경}
\newcommand{\dd}{\mathrm{d}}
\newcommand{\rxn}{\mathrm{rxn}}
\newcommand{\eq}{\mathrm{eq}}
\newcommand{\rev}{\mathrm{rev}}
\newcommand{\irr}{\mathrm{irr}}
\newcommand{\oc}{\mathrm{oc}}
\newcommand{\config}{\mathrm{config}}
\newcommand{\vib}{\mathrm{vib}}
\newcommand{\eff}{\mathrm{eff}}
\newcommand{\app}{\mathrm{app}}
\newcommand{\hys}{\mathrm{hys}}
\newcommand{\cell}{\mathrm{cell}}
\newcommand{\kB}{k_{\mathrm B}}
\newcommand{\code}[1]{\texttt{#1}}
\newcommand{\avg}[1]{\langle #1\rangle}
\renewcommand{\theequation}{2.\arabic{equation}}
\renewcommand{\thesection}{2.\arabic{section}}
